% ============================================================
% Section 1.3: Project Engineering Objectives
% ============================================================
\section{Project Engineering Objectives}
\label{sec:project_engineering_objectives}

The primary objective of this project is to develop an integrated, cost-effective, non-invasive HMI that provides individuals with locked-in syndrome (LIS) an assistive interface for communication and navigation. To attain this overall goal, the following specific engineering objectives were established:

\begin{enumerate}[label=\textbf{\arabic*.}]
    \item \textbf{Bio-Potential Acquisition Hardware Design:} To design and assemble a high-gain, low-noise analog front-end circuit capable of reliably capturing microvolt-level corneo-retinal potentials while rejecting power-line interference and environmental electromagnetic noise.
    \item \textbf{Real-Time Signal Processing Framework:} To implement real-time digital filtering and thresholding algorithms within a hybrid MATLAB-Arduino processing architecture to extract and classify the temporal characteristics of voluntary eye blinks.
    \item \textbf{Finite State Machine Control Integration:} To engineer a deterministic Finite State Machine (FSM) that translates multi-blink sequences into operational commands, allowing a single biopotential channel to manage two distinct subsystems while enforcing safety protocols.
    \item \textbf{Virtual Optical Speller Development:} To construct an auto-scanning virtual keyboard featuring a Graphical User Interface (GUI) configured to reduce cognitive load, enabling users to compose, edit, and clear text.
    \item \textbf{3D-Printed Prototype Wheelchair Implementation:} To construct a functional 3D-printed motorized wheelchair prototype that receives wireless directional commands, serving as a scaled experimental model to validate wheelchair navigation and obstacle-avoidance logic.
\end{enumerate}
